\paragraph{Expanded learning-rate screen.}
Because the original development selection lay at the upper edge of its rate grid, we performed a separately specified post hoc screen using the original three development seeds (201--203). Each of the four credit rules received the same six-rate search on both the aligned and opposed tasks: $\{0.003,0.01,0.03,0.1,0.3,1\}$ for Adam and $\{0.03,0.1,0.3,1,3,10\}$ for SGD. We reran the original rates alongside the added rates, giving every fit 16,384 updates and selecting its minimum validation NMSE among checkpoints at 0, 64, 256, 1,024, 2,048, 4,096, 8,192, 12,288 and 16,384 updates. All other task, initialization, optimizer and conductance-bound settings remained those of the original opponent study. The screen comprises 288 retained fits; no test outcome entered rate selection.

Before inspecting these outcomes, we specified that an improvement of at least 0.01 NMSE and at least 10\% in mean opposed-task validation error for either broadcast rule would trigger a follow-up on the 20 previously used fresh seed blocks. The first criterion compared the best added rate with the best original-grid rate at the same update budget; a pre-outcome addendum also compared the best rate across the whole grid with the originally selected rate (0.03 for Adam and 0.3 for SGD). Any triggered follow-up would retune all four rules on both tasks using development validation alone and would be reported as reuse of existing blocks. Neither criterion was met. No added rate improved the opposed-task broadcast minimum; within the original grid, reducing Adam's rate to 0.003 lowered mean validation NMSE by approximately 0.00647 (0.89\%), and reducing the unit-broadcast SGD rate to 0.1 lowered it by 0.03348 (4.46\%). Calibrated-broadcast SGD retained rate 0.3. We therefore performed no additional fresh-block fits. These results address a finite grid at the stated update budget and do not establish optimizer-independent convergence.
